%---------------------------------------------------------------------
% Title page and a single orientation page.
%---------------------------------------------------------------------
\thispagestyle{empty}

% Band depth measured from the rendered page: Nastaliq at Scale=1.45 sets
% taller lines than the nominal point sizes suggest.
\begin{tikzpicture}[remember picture, overlay]
  \fill[bmblued] (current page.north west) rectangle
        ([yshift=-12.5cm]current page.north east);
  \fill[bmaccent] ([yshift=-12.5cm]current page.north west) rectangle
        ([yshift=-12.7cm]current page.north east);
\end{tikzpicture}

\vspace*{0.9cm}
{\color{white}
  {\large علامہ اقبال اوپن یونیورسٹی، اسلام آباد}\par
  \vspace{3pt}
  {شعبۂ کمپیوٹر سائنس}\par
  \vspace{18pt}
  {\Huge\bfseries بنیادی معلوماتی و ابلاغی ٹیکنالوجی}\par
  \vspace{5pt}
  {\large\bfseries \textenglish{Basics of Information and Communication Technology}}\par
  \vspace{10pt}
  {\LARGE\bfseries اہم ترین سوالات --- مختصر ایڈیشن}\par
  \vspace{12pt}
  {\large کورس کوڈ \textenglish{1431} \ (بعض اوقات \textenglish{5403})}\par
  \vspace{4pt}
  % \par inside the group: colour is applied when the paragraph breaks, so
  % closing the group first would revert to black ink on the dark band.
  {نو یونٹ \ \textbullet\ \ سولہ سوالات \ \textbullet\ \ بی اے، اے ڈی ای، بی ایڈ، بی ایس\par}
}

\vspace{1.0cm}

\begin{tcolorbox}[enhanced, colback=bmtint, colframe=bmblue, boxrule=0.6pt,
                  arc=3pt, left=12pt, right=12pt, top=10pt, bottom=10pt]
\textbf{\color{bmblue}یہ کتابچہ کیا ہے؟}\ \
اے آئی او یو کی سرکاری درسی کتاب \textenglish{\emph{Basics of Information and
Communication Technology}} (کورس کوڈ \textenglish{1431/5403}، شعبۂ کمپیوٹر
سائنس، پہلا ایڈیشن \textenglish{2013}) کے ہر یونٹ کے آخر میں دیے گئے
\textenglish{Self-Assessment Questions} میں سے \textbf{سولہ سب سے اہم} سوالات
کا مختصر حل۔ یونیورسٹی اسائنمنٹ اور فائنل پرچہ اسی ذخیرے سے بناتی ہے۔

\medskip
\textbf{\color{bmblue}اصطلاحات کے بارے میں۔}\ \
اس مضمون کی تکنیکی اصطلاحات انگریزی ہی میں لکھی گئی ہیں، کیونکہ پرچے میں بھی
انہی کی توقع ہوتی ہے۔ ہر اصطلاح کی اردو وضاحت ساتھ دی گئی ہے، اور آخر میں
\hyperref[sec:lughat]{اصطلاحات کی فہرست} بھی موجود ہے۔

\medskip
\textbf{\color{bmblue}حیثیت۔}\ \
ذاتی مطالعے کے لیے۔ یہ اے آئی او یو کی مطبوعہ کتاب نہیں اور نہ ہی اسے یونیورسٹی
کی توثیق حاصل ہے۔ اسے من و عن نقل کر کے اسائنمنٹ میں جمع کرانا منع ہے۔
\end{tcolorbox}

\clearpage

%---------------------------------------------------------------------
\section*{پہلے یہ پڑھ لیجیے}
\markboth{پہلے یہ پڑھ لیجیے}{}

سوال پڑھیے۔ جواب کو کاغذ سے ڈھانپیے۔ پہلے خود کاپی پر لکھیے، \emph{پھر} ملائیے۔
اس مضمون میں زیادہ تر سوال ''تعریف کیجیے''، ''فرق واضح کیجیے'' اور ''مثالوں سے
سمجھائیے'' کی قسم کے ہوتے ہیں، اس لیے تیاری کا بہترین طریقہ یہ ہے کہ ہر تصور کی
\textbf{ایک سطری تعریف} اور \textbf{دو مثالیں} یاد ہوں۔

\medskip
\begin{itemize}\setlength{\itemsep}{2pt}
  \item \textbf{سبز خانہ} --- مختصر جواب، جو کم از کم لکھنا ضروری ہے۔
  \item \textbf{نارنجی خانہ} --- وہ بنیادی اصول یا فہرست جو ایک سے زیادہ سوالوں
        میں کام آتی ہے۔
  \item \textbf{سرخ خانہ} --- اس قسم کے سوال کی مخصوص غلطی۔
\end{itemize}

\begin{nukta}[پرچے میں جواب لکھنے کا سانچہ]
''فرق واضح کیجیے'' والے سوال میں \textbf{جدول} بنائیے --- ممتحن کو نکات فوراً نظر
آتے ہیں اور نمبر پورے ملتے ہیں۔ ''تعریف کیجیے'' والے سوال میں ترتیب یہ رکھیے:
\textbf{(۱)} ایک سطری تعریف، \textbf{(۲)} دو یا تین اہم خصوصیات،
\textbf{(۳)} کم از کم دو مثالیں۔ صرف تعریف لکھ دینے سے آدھے نمبر ملتے ہیں۔
\end{nukta}

\medskip
\noindent
\textbf{\color{bmblue}پرچے میں کس یونٹ کا کتنا وزن ہے}

\medskip
\begin{center}
\begin{tabular}{@{}p{5.0cm} p{1.8cm} p{1.6cm} p{6.0cm}@{}}
\toprule
\textbf{یونٹ اور موضوع} & \textbf{وزن} & \textbf{سوالات} & \textbf{اصل امتحان کس چیز کا ہے} \\
\midrule
\textenglish{1} \textenglish{ICT} کا تعارف        & درمیانہ & \textenglish{2} & فوائد و نقصانات؛ عملی اطلاقات \\
\textenglish{2} کمپیوٹر کی ساخت                    & زیادہ   & \textenglish{3} & \textenglish{CPU}، میموری، نسلیں \\
\textenglish{3} اِن پُٹ آلات                        & زیادہ   & \textenglish{2} & \textenglish{OCR/OMR/MICR} کا فرق \\
\textenglish{4} آؤٹ پُٹ آلات                       & زیادہ   & \textenglish{2} & پرنٹر کی اقسام کا موازنہ \\
\textenglish{5} سافٹ ویئر                          & زیادہ   & \textenglish{2} & سسٹم بمقابلہ ایپلیکیشن \\
\textenglish{6} آپریٹنگ سسٹم                       & درمیانہ & \textenglish{1} & افعال کی فہرست \\
\textenglish{7} ڈیٹا کمیونیکیشن و نیٹ ورکنگ        & زیادہ   & \textenglish{2} & \textenglish{LAN/WAN}، ٹوپولوجی \\
\textenglish{8} ملٹی میڈیا                          & درمیانہ & \textenglish{1} & اجزا اور اطلاقات \\
\textenglish{9} کمپیوٹر زبانیں                      & درمیانہ & \textenglish{1} & \textenglish{Compiler} بمقابلہ \textenglish{Interpreter} \\
\bottomrule
\end{tabular}
\end{center}

\medskip
\noindent
یونٹ \textenglish{2}، \textenglish{3}، \textenglish{4}، \textenglish{5} اور
\textenglish{7} مل کر زیادہ تر نمبر اٹھاتے ہیں۔ اگر وقت کم ہو تو انہی سے شروع
کیجیے۔
